\section{Conclusion}
\label{sec:conclusion}
We have presented GNN-QMIX, a value-decomposition method for cooperative conflict
resolution whose parameters do not depend on the number of vehicles. Following the
entity-graph principle of TransfQMix, observations and the state are entity matrices. A
GCN shared by all vehicles encodes each vehicle's sensor-radius observation graph into an
embedding and its individual action values, and a GAT over the agent graph generates the
weights of a monotonic mixer from those embeddings and the global state. The mixer
satisfies the QMIX monotonicity condition at every fleet size and is used only in training;
at execution each vehicle acts greedily on its own action values. The same parameters can
therefore be trained jointly at several fleet sizes and executed at others. A second stage
places the learned policy in a hybrid controller, in which a PID waypoint-tracking law
issues the heading and speed commands and the learned policy overrides it for a vehicle
only while another vehicle is inside that vehicle's sensor radius.

The evaluation protocol uses the BlueSky simulator with joint training at $n\in\{3,6,9\}$,
zero-shot evaluation at $n\in\{10,20,50\}$, and QMIX, VDN, and IQL as baselines. The results
are not yet available; the expectations stated in this paper are hypotheses.
% TODO(results): replace the previous sentence with a summary of the measured outcomes
% (summary_training.csv, summary_zero_shot.csv) and state which hypotheses of
% Section~\ref{subsec:hypotheses} were supported.

\textbf{Limitations and future work.} The present method and experiments leave out several
elements of the original problem statement. (i) Communication: there is no broadcast
channel, no communication latency, and no recurrent belief state; each vehicle acts on its
current sensor data only, and the characterization of performance against latency is
future work.
% TODO(authors): deferred per decision 2026-11-25 (GAP-07/GAP-09/GAP-12) -- latency \delta,
% broadcast channel, GRU belief, and theoretical claim 3.
(ii) Non-cooperative traffic and obstacles: every aircraft in the experiments is a
cooperating agent, and there are no static obstacles.
% TODO(authors): deferred per decision 2026-11-25 (GAP-07/GAP-09/GAP-12) -- intruders and
% static obstacles.
(iii) Actions: the learned policy selects among 15 discrete heading and speed changes;
continuous actions, and the continuous-action baselines they would permit, are not
considered.
% TODO(authors): deferred per decision 2026-11-25 (GAP-07/GAP-09/GAP-12) -- continuous
% actions, MAPPO and GCBF baselines.
(iv) Vehicle model: the experiments use a point-mass transport aircraft model at a single
flight level with a 5~NM separation minimum. The method is motivated by UAS and AAM
operations, but no UAS dynamics, separation standards, or vertical maneuvers are evaluated.
% TODO(authors): confirm UAS vs ATM framing (Q-01).
(v) Density: the zero-shot tests keep the sector size fixed, so fleet size and traffic
density change together and their effects are not separated. (vi) Safety: the method does
not provide a formal safety certificate, and the activation rule of the hybrid controller
is a distance threshold. A possible extension is to replace the radius trigger by a learned
safety score computed from each vehicle's node embedding, with a threshold calibrated
against a control barrier function condition such as that of GCBF~\cite{gcbf}, which would
connect the switching logic to certificate-based control.
% NOTE (confirmed by the authors 2026-11-25, D-031): that the learned safety score s^i / CBF-calibrated threshold \theta
% should be presented as a possible extension (it was the switching rule of the extended
% abstract).
(vii) The details of training in the hybrid stage, the PID law, and the definition of a
per-vehicle mission-success metric for this environment remain to be fixed.
% TODO(authors): mission-success definition (GAP-18).
